\section{Application: policy-topic coding}\label{sec:application}
% TODO-AUTHOR: 5.1 and the first paragraph of 5.2 are Claude's factual drafts from data/README.md and
% the annotation logs; the "choice of models" paragraph is the author's text (manuscript/model_and_environment.md).
% 5.3 waits for the analysis (results/application/summary.md).

\subsection{Data}\label{sec:app_data}

We use bill titles from the Comparative Agendas Project's Congressional Bills dataset (version 19.3),
in which every bill carries one major policy topic assigned by trained human coders under the
Policy Agendas codebook. From the 37,007 bills of the 111th to 116th Congresses (2009--2020) with a
title and one of the 20 major topics we drew a stratified sample of 800: 240 bills coded as Health
(topic 3; 11.6\% of the population) and 560 from the other 19 topics in proportion to their
frequency, with a fixed seed. The binary target is whether the bill's major topic is Health. A
random 200 bills form a validation set whose human codes are never used in estimation; the
remaining 600 are the working set from which 30 anchors (5\%) are drawn for the anchored fit. Titles
have a median length of 25 words. The derived file (bill identifier, Congress, title, CAP major
topic and our labels) is redistributed under the dataset's CC BY-NC-SA 4.0 licence with attribution.

\subsection{Annotation design}\label{sec:app_design}

Each bill was labelled by $M = 3$ model families $\times$ $P = 5$ prompts $\times$ $R = 3$ decoding
runs at temperature 0.7, giving 45 labels per bill and 12,000 calls per provider, plus one
temperature-0 pass of prompt 1 (800 calls per provider). The five prompt templates, written by the
authors and reproduced verbatim in the supplementary material, range from a bare question to a
codebook definition, a contrast with other topics, a role instruction and an instruction to
consider the purpose before answering. Every call asked for a single digit with a 16-token output
budget; the first digit of the answer was parsed as the label, an answer without a 0/1 digit was
re-queried once and then recorded as missing. Calls were made from a script that stores every raw
response with timestamp, model version, prompt, run and temperature (Table~\ref{tab:providers}).

The three model families were chosen to span the design space that the method is meant to
describe rather than to maximise accuracy. First, each family comes from a different developer
with a different training corpus and alignment procedure (Anthropic, Google, Alibaba), so that the
item-by-model effect $\phi_{im}$ captures genuine between-family dependence rather than variation
among versions of one model. Second, the two commercial models are the efficiency tier of their
families (Claude Haiku 4.5, Gemini 3.5 Flash), which is the tier practitioners use for annotation
at scale: \citet{Carlson2025} report that such models match or approach the larger tiers on binary
classification at a fraction of the cost, and cost is what makes the model $\times$ prompt
$\times$ run design affordable. Third, one family is an open-weights model run locally (Qwen 2.5
7B Instruct, 4-bit quantised, served by Ollama): its weights and digest are archived, so this arm
of the ensemble is exactly reproducible, which no API-served model can guarantee once its version
is retired \citep{Barrie2025}. Larger or reasoning-oriented models were deliberately excluded:
they raise cost by one to two orders of magnitude, their suitability for plain classification is
unclear \citep{Carlson2025}, and the question of the paper is not which model is best but how
dependent the members of an ensemble are. Model identifiers, quantisation, decoding parameters,
query dates, call counts, missing rates and cost are given in Table~\ref{tab:providers}; all raw
responses are archived.
\begin{table}[htbp]
\centering
\footnotesize
\setlength{\tabcolsep}{4pt}
\caption{LLM configurations used in the application. Calls count the $800 \times 5 \times 3$ labels at temperature 0.7 plus the 800-call temperature-0 pass of prompt 1; missing = answers without a 0/1 digit after one retry; cost at list prices from the logged token counts (Haiku 4.5: USD 1.00 / 5.00 per million input / output tokens; Gemini 3.5 Flash: USD 1.50 / 9.00)}\label{tab:providers}
\begin{tabular}{llll}
\toprule
Provider & Anthropic & Google & Alibaba (open weights) \\
\midrule
Model identifier & claude-haiku-4-5-20251001 & gemini-3.5-flash & qwen2.5:7b-instruct \\
Version / digest & as served (API string) & as served, thinking off & digest 845dbda0ea48 \\
Quantisation & -- & -- & Q4\_K\_M, 7.6B \\
Endpoint & Messages API & Gemini API & Ollama 0.34.4 (local) \\
Temperature & 0.7 (and 0) & 0.7 (and 0) & 0.7 (and 0) \\
Max output tokens & 16 & 16 & 16 \\
Query dates & 25 Sept 2026 & 25--26 Sept 2026 & 25 Sept 2026 \\
Calls & 12,800 & 12,800 & 12,800 \\
Missing (\%) & 0.00 & 0.08 & 0.00 \\
Cost (USD) & 1.63 & 1.68 & 0 \\
\bottomrule
\end{tabular}
\end{table}


\subsection{Results}\label{sec:app_results}
% TODO-AUTHOR: to be written from results/application/summary.md (agreement, variance shares,
% evaluation against CAP, per-configuration accuracy, calibration Fig. 5, posterior predictive check Fig. 6).
[To be completed when the analysis has run.]
